\documentclass[11pt]{article}
\usepackage{mathblatt}
% Präambel des Setzers (werkzeuge/setzer-praeambel.tex, wortgleich); dazu \gteil (Muster T2U)
\makeatletter
\setlength{\mbpfbe}{0mm}% keine Punkte (pruefung.md § 5: nur beim Prüfstein)
\newlength{\gEin}\setlength{\gEin}{0pt}
\let\gSchrift\relax
\renewcommand{\pfkreuzzeile}[1]{\par\noindent\hangindent3.2em\hangafter1\hspace*{1.6em}$\square$\ #1\par}
\newcommand{\gkreuz}[1]{$\square$\ #1\hspace{2em}}
\newcommand{\gkreuzzeile}[2]{\par\noindent\hangindent2.2em\hangafter1\mbox{#1\,$\square$}\ #2\par}% Buchstabe, Option: eine je Zeile (neben Skizze)
% links, was man ansieht, rechts, was man tut; Spaltengrenze überall an derselben Stelle (gemeinsam.md „Satz“)
% v0.17: Fußeintrag „Nr: Wert“ – kurz in einer Box (nie getrennt), lang umbrechbar mit der Nummer am Anfang
% (vorher \mbox: lange Lösungen liefen über den Rand, AS40 138 pt); robust, weil \fusshilfe in die .aux schreibt
\NewDocumentCommand{\gfuss}{m m}{\sbox0{#1:~#2}\ifdim\wd0<0.45\textwidth\mbox{#1:~#2}\else\mbox{#1:}\nobreak\ #2\fi}
\newsavebox{\gbox}\newlength{\gLinks}
% v0.10 (Lehrer 13.10.): keine starre Halbierung – die linke Spalte ist so breit wie die Figur (+ 6 mm), die
% Antwort steht direkt daneben; ist die Figur breiter als 0,62 Zeile, steht die Antwort darunter
\newcommand{\gzwei}[2]{\par\smallskip\noindent\sbox{\gbox}{#1}%
  \setlength{\gLinks}{\dimexpr\wd\gbox+6mm\relax}%
  \ifdim\gLinks>0.62\linewidth
    \usebox{\gbox}\par\smallskip\noindent\begin{minipage}{\linewidth}#2\end{minipage}%
  \else
    \begin{minipage}[t]{\gLinks}\vspace{0pt}\usebox{\gbox}\end{minipage}%
    \begin{minipage}[t]{\dimexpr\linewidth-\gLinks\relax}\vspace{0pt}\raggedright #2\end{minipage}%
  \fi\par}
% 16.10. (Lehrer, PA10 S. 1): Text links, Koordinatensystem rechts daneben; ist das System breiter als 0,58 Zeile,
% steht es wie bisher unter dem Text
% Passt das System knapp nicht mehr auf die Seite (bis 25 % zu hoch), wird es auf den Rest verkleinert statt
% die Seite halb leer zu lassen
\newlength{\gRest}\newlength{\gHoch}
\newcommand{\gtextks}[2]{\sbox{\gbox}{#1}%
  \setlength{\gHoch}{\dimexpr\ht\gbox+\dp\gbox\relax}%
  \ifdim\pagegoal<\maxdimen\setlength{\gRest}{\dimexpr\pagegoal-\pagetotal-4\baselineskip\relax}%
    \ifdim\gHoch>\gRest\ifdim\gRest>0.75\gHoch
      \sbox{\gbox}{\resizebox*{!}{\gRest}{\usebox{\gbox}}}\fi\fi\fi
  \ifdim\wd\gbox>0.58\linewidth #2\par\smallskip\noindent\gpassend{\usebox{\gbox}}\par
  \else\begin{minipage}[t]{\dimexpr\linewidth-\wd\gbox-5mm\relax}\raggedright #2\end{minipage}\hfill
    \begin{minipage}[t]{\wd\gbox}\vspace{0pt}\usebox{\gbox}\end{minipage}\par\fi}
% v0.9 (12.10., pruefung.md § 5 „Marke mit B, Bedingungen rechts“): drittes Argument von pfaufg = Bedingungen
% der Prüfung (v0.10: \pfohnetr = TR durchgestrichen, „FOR“), klein grau im rechten Rand auf der Grundlinie der Nummer, Zeilen mit \\
\newcommand{\gbedrand}[2]{\ifblank{#2}{}{\rlap{\hspace*{\dimexpr#1+2mm\relax}%
  \parbox[t]{15mm}{\raggedright\scriptsize\leavevmode\color{mbgrau}\strut#2}}}}
% Teil eines Bündels (a, b …): Marke links im Rand, Bedingungen rechts im Rand, auf der Zeile des Teils
\newcommand{\gteilrand}[2]{\llap{\mbpfrandmarke{#1}\hspace*{\mbpfnr}}\gbedrand{\linewidth}{#2}}
% Teil eines Bündels: Buchstabe in eigener Spalte, Inhalt eingerückt (\gEin = Einzug, damit die Spaltengrenze von
% \gzwei an derselben Stelle bleibt)
\newenvironment{gbteil}[3]{\par\smallskip\noindent\gteilrand{#2}{#3}\makebox[\mbpfteil][l]{\textbf{#1}}%
  \begin{minipage}[t]{\dimexpr\linewidth-\mbpfteil\relax}\setlength{\gEin}{\mbpfteil}\raggedright\strut\ignorespaces}%
 {\end{minipage}\par\addvspace{6pt}}% v0.18 (PG80 Nr. 1): Abstand, damit die Linie von a) nicht an b) rutscht
% Grafik höchstens so breit wie die Spalte (Bündel nebeneinander, Koordinatensystem ohne Karo)
\newsavebox{\gpbox}
\newcommand{\gpassend}[1]{\sbox{\gpbox}{#1}\ifdim\wd\gpbox>\linewidth\typeout{GPASSEND: verkleinert \the\wd\gpbox\space auf \the\linewidth}\resizebox{\linewidth}{!}{\usebox{\gpbox}}\else\usebox{\gpbox}\fi}
% Skizzenfeld (v0.10, Lehrer 13.10.): Karo ab dem linken Rand über die Zeile, „Skizze“ klein grau oben links im Feld
% (wie \pfskizzenfeld der Vorlage); der rechte Rand bleibt den Bedingungen
\newcommand{\gskizze}[1]{\par\smallskip\noindent\begin{tikzpicture}
  \clip (0,0) rectangle ({\linewidth-1mm},{#1*10mm});
  \draw[black!16,line width=0.3pt,step=5mm] (0,0) grid ({\linewidth},{#1*10mm});
  \node[anchor=north west,inner sep=2pt,font=\scriptsize,text=mbgrau,fill=white] at (0,{#1*10mm}) {Skizze};
  \end{tikzpicture}\par}
\RenewDocumentEnvironment{pfaufg}{m m m}{%
  \begin{lrbox}{\mbaufgabebox}\begin{minipage}[t]{\linewidth}%
  \gSchrift\noindent\hspace*{\gEin}\mbpfrandmarke{#1}\gbedrand{\linewidth-\gEin}{#3}%
  \begin{minipage}[t]{\mbpfnr}\strut\textbf{#2}\end{minipage}%
  \begin{minipage}[t]{\dimexpr\linewidth-\gEin-\mbpfnr\relax}\raggedright\strut\ignorespaces}%
 {\end{minipage}%
  \end{minipage}\end{lrbox}%
  \setlength{\mbaufgabehoehe}{\dimexpr\ht\mbaufgabebox+\dp\mbaufgabebox\relax}%
  \ifdefempty{\gkopfnach}{\par\addvspace{5pt}\Needspace*{\mbaufgabehoehe}}{%
    % Stufenkopf, der mit dieser (ersten) Aufgabe zusammenbleibt: Platz für beide verlangen, dann den Kopf setzen
    % (12.10., Quadratische: ein Kopf stand allein unten, die Aufgabe mit Koordinatensystem auf der nächsten Seite)
    % v0.17 (AS40 „Terme deuten“ allein unten): hinter dem Kopf steht das Sprungziel (\special, ein Knoten, nach
    % dem die Seite brechen durfte) – darum \nobreak nach dem Kopf; Vorhalt 56 pt = Kopf mit Linie und Abständen
    \par\Needspace*{\dimexpr\mbaufgabehoehe+56pt\relax}\gkopfnach\par\nobreak\gdef\gkopfnach{}\addvspace{5pt}}%
  \noindent\usebox{\mbaufgabebox}\par\addvspace{5pt}}
% v0.18 (plan.md § 10, 15.10. l): Abschnittsbuchstabe A, B, C vor dem Namen im Kopf – als Vorsatz, damit der
% Name im TeX gleich bleibt (Bestellkarte, Sprungziele lesen ihn)
\def\gbstakt{}\newcommand{\gbstvor}[1]{\gdef\gbstakt{#1}}
\NewCommandCopy\gpfstufekopfalt\pfstufekopf
\RenewDocumentCommand{\pfstufekopf}{o m m}{\IfValueTF{#1}{\gpfstufekopfalt[#1]}{\gpfstufekopfalt}%
  {\ifx\gbstakt\empty\else\gbstakt\enspace\fi #2}{#3}\gdef\gbstakt{}}
% v0.18 (16.10. e): Begriff grau, ohne Nummer, ohne Rahmen über der ersten Aufgabe der Stufe
\newcommand{\gbegriff}[1]{\par\nobreak\noindent\hspace*{\mbpfnr}\parbox[t]{\dimexpr\linewidth-\mbpfnr\relax}{\raggedright\small\color{mbgrau}#1}\par\nobreak\addvspace{3pt}}
% v0.18 (16.10. g): Lösungsblatt – Kopf mit kräftigem Strich, Abschnitt klein grau als Zwischenzeile
\newcommand{\glkopf}[1]{\noindent{\large\bfseries Lösungen \textperiodcentered{} #1}\par\vspace{3pt}\hrule height 1.2pt\par\medskip}
\newcommand{\glzwischen}[1]{\par\addvspace{8pt}\Needspace*{4\baselineskip}\noindent{\footnotesize\color{mbgrau}#1}\par\nobreak\vspace{1pt}}
% v0.18 (15.10. o, w): Blattende – Kern ohne Rahmen, dann graue Zeile „Mehr: n Aufgaben“
\newcommand{\gkern}[1]{\par\addvspace{14pt}\Needspace*{6\baselineskip}\noindent{\bfseries Kern}\par\nobreak\smallskip
  {\small\raggedright #1\par}}
\newcommand{\gmehr}[1]{\par\addvspace{10pt}\noindent{\footnotesize\color{mbgrau}Mehr: #1 Aufgaben}\par}
% Stufenkopf vormerken: \gkopfvor{<Kopf>} – die nächste pfaufg setzt ihn, sobald beide zusammen Platz haben
\def\gkopfnach{}
\newcommand{\gkopfvor}[1]{\gdef\gkopfnach{#1}}
% Teilaufgabe mit Skizze links in gzwei: Buchstabe in eigener Spalte, darunter Text und Skizze
\newcommand{\gteil}[2]{\makebox[\mbpfteil][l]{\textbf{#1}}\parbox[t]{\dimexpr0.5\textwidth-\gEin-\mbpfnr-5mm-\mbpfteil\relax}{\raggedright #2}}
% Stufenkopf der Vorlage (mathblatt.sty 2026-10-12: Linie, Abstand, 13 pt, \pfjahre) ohne ihren Vorhalt von
% 14 Zeilen: der Setzer fordert vor dem Kopf 0,16 Seitenhöhe (Kopf + erste Aufgabe), mehr nicht
\setlength{\mbpfstufevorhalt}{0pt}
\makeatother
\newcommand{\karo}[1]{\par\smallskip\noindent\begin{tikzpicture}
  \clip (0,0) rectangle ({\linewidth-1mm},{#1*10mm});
  \draw[black!16,line width=0.3pt,step=5mm] (0,0) grid ({\linewidth},{#1*10mm});
  \end{tikzpicture}\par}
\newcommand{\kk}[1]{$\square$\,#1\hspace{0.9em}}
\newcommand{\linie}[1]{\,{\color{black!45}\rule[-1pt]{#1}{0.4pt}}\,}
\newcommand{\pkt}{\textperiodcentered{}}
% v0.14 (14.10.): Stufenkopf mit fünf Punkten (Option = Jahre für die Bestellkarte) oder „zuletzt 19“
\NewDocumentCommand{\gpunkte}{o m}{#2}
\newcommand{\gzuletzt}[1]{zuletzt #1}
% v0.7 (12.10.): Schülerfrage unter dem Stufenkopf; Ankreuzoption als Bild (Buchstabe und Kästchen darunter)
\newcommand{\gfrage}[1]{\par\nobreak\noindent\hspace*{\mbpfnr}{\small\itshape #1}\par\nobreak\addvspace{2pt}}
\newlength{\gbildbreite}\setlength{\gbildbreite}{0.32\linewidth}% Setzer setzt sie je Reihe (vier Bilder: 0,24)
\newcommand{\gbildoption}[3]{\begin{minipage}[t][#3][b]{\gbildbreite}\centering #2\par\smallskip\mbox{#1\,$\square$}\end{minipage}}
% Lineare (12.10.): Bild mit Buchstabe ohne Kästchen (Diagramm zuordnen); Aussage mit „wahr / falsch“ zum Ankreuzen
\newcommand{\gbildname}[3]{\begin{minipage}[t][#3][b]{\gbildbreite}\centering #2\par\smallskip\mbox{\textbf{#1}}\end{minipage}}
\newcommand{\gwahrfalsch}[2]{\par\noindent\hangindent2.2em\hangafter1\makebox[2.2em][l]{(#1)}#2\hfill\mbox{wahr $\square$\quad falsch $\square$}\par}

\newcommand{\gunter}[2]{\par\addvspace{6pt}\noindent{\bfseries #1\enspace #2}\par\nobreak\vspace{2pt}}
\newcommand{\gsternzeile}{\par\vspace{-4pt}\noindent{\footnotesize\color{mbgrau}* nur FOR}\par\medskip}
% 17.10.: Skizze oben bündig mit der ersten Textzeile (beide Spalten hängen an der Oberkante); der Rechenraum (Karo)
% beginnt direkt unter dem Text, links neben dem Graph, und läuft unter dem Graph über die ganze Zeile weiter
% (Streifen links = auf 5 mm aufgerundete Höhe des Graphen unter dem Text, Raster durchgehend; die Fläche des Streifens
% wird von der Karohöhe unter dem Graph abgezogen, auf 0,5 cm aufgerundet – die Gesamtfläche Rechenraum bleibt mindestens #3 × Zeile); #3 = Zeilen Karo in cm
% unter dem Graph (0 = kein Karo). Kein Verkleinern des Graphen.
\newsavebox{\gtbox}\newlength{\gWl}\newlength{\gDh}
\newcommand{\gkaroohne}[1]{\nointerlineskip\noindent\begin{tikzpicture}
  \clip (0,0) rectangle ({\linewidth-1mm},{#1*10mm});
  \draw[black!16,line width=0.3pt,step=5mm] (0,0) grid ({\linewidth},{#1*10mm});
  \end{tikzpicture}\par}
\newcommand{\gtextkaro}[3]{\sbox{\gbox}{#1}%
  \ifdim\wd\gbox>0.58\linewidth #2\par\smallskip\noindent\gpassend{\usebox{\gbox}}\par\ifnum#3>0 \karo{#3}\fi
  \else
    \setlength{\gWl}{\dimexpr\linewidth-\wd\gbox-5mm\relax}%
    \sbox{\gtbox}{\begin{minipage}[t]{\gWl}\raggedright #2\par\end{minipage}}%
    \setlength{\gDh}{\dimexpr\ht\gbox+\dp\gbox-\ht\gtbox-\dp\gtbox-6pt\relax}%
    \leavevmode
    \ifnum#3>0\ifdim\gDh>0pt
      \pgfmathsetmacro{\gxs}{floor(\gWl/5mm)*5mm}\pgfmathsetmacro{\gys}{ceil(\gDh/5mm)*5mm}%
      \begin{minipage}[t]{\gWl}\usebox{\gtbox}\par\vspace{6pt}\nointerlineskip\noindent
        \begin{tikzpicture}
        \clip (0,0) rectangle (\gxs pt,\gys pt);
        \draw[black!16,line width=0.3pt,step=5mm] (0,0) grid (\gxs pt,\gys pt);
        \end{tikzpicture}\par\end{minipage}%
    \else\usebox{\gtbox}\fi\else\usebox{\gtbox}\fi
    \hfill\begin{minipage}[t]{\wd\gbox}\vspace{-\dimexpr\ht\strutbox+2.4ex\relax}\usebox{\gbox}\end{minipage}\par
    \ifnum#3>0 \ifdim\gDh>0pt
      \pgfmathsetmacro{\gnn}{max(0,ceil(2*(#3-(\gys/28.4528)*(\gxs/\linewidth)))/2)}%
      \ifdim\gnn pt>0pt\gkaroohne{\gnn}\fi
    \else\karo{#3}\fi\fi
  \fi}
\begin{document}

\pfheftstil
\pflaufkopf[P10 \pkt{} Prüfung]{Trigonometrie}{PG60}\pfseitevon
\pfheftkopf{Trigonometrie}{}
\pfstufenzeile{Berechnungen im rechtwinkligen Dreieck \pkt{} Berechnungen in beliebigen Dreiecken}
\gsternzeile

% ===================== A Berechnungen im rechtwinkligen Dreieck =====================
\gkopfvor{\gbstvor{A}\pfstufekopf{Berechnungen im rechtwinkligen Dreieck}{}\gunter{A1}{Winkel gesucht}}
% DRUCK 2026-FOR-K4b deckt 2026-FOR-K4b 2026-EBR-K3b
\begin{pfaufg}{nach~P10\\’26}{1.}{}
\gtextkaro{\begin{tikzpicture}[baseline=(current bounding box.north),line width=0.7pt,font=\small]\draw (0,0) -- (4.6,0) -- (5.9,2.92) -- (1.3,2.92) -- cycle;\draw[densely dotted] (4.6,0) -- (5.9,0) -- (5.9,2.92);\pfrw[0.25]{(5.9,0)}{90}\node[anchor=north east,inner sep=1.5pt] at (0,0) {$A$};\node[anchor=north east,inner sep=1.5pt] at (4.6,0) {$B$};\node[anchor=north west,inner sep=1.5pt] at (5.9,0) {$F$};\node[anchor=south west,inner sep=1.5pt] at (5.9,2.92) {$C$};\node[anchor=south east,inner sep=1.5pt] at (1.3,2.92) {$D$};\node[anchor=north,inner sep=1.5pt,font=\footnotesize] at (5.25,0) {$13$ cm};\node[anchor=west,inner sep=1.5pt] at (5.9,1.46) {$h$};\node[rotate=66,anchor=south,inner sep=2pt,font=\footnotesize] at (5.25,1.46) {$32$ cm};\pfwinkel[0.5][0.65]{(4.6,0)}{0}{66}{\varepsilon}\end{tikzpicture}}{Bei dem Viereck $ABCD$ sind die gegenüberliegenden Seiten jeweils parallel zueinander. Weise nach, dass die Größe des Winkels $\varepsilon$ etwa $66^\circ$ beträgt.}{4}
\fusshilfe{\gfuss{1}{66,0°}}
\end{pfaufg}
% DRUCK 2024-OS-K6b deckt 2024-OS-K6b
\begin{pfaufg}{nach~P10\\’24}{2.}{}
\gtextkaro{\begin{tikzpicture}[baseline=(current bounding box.north),line width=0.7pt,font=\small]\draw (0,2.297) -- (0,0) -- (2.04,0) -- cycle;\pfrw[0.25]{(0,0)}{0}\node[anchor=south,inner sep=1.5pt] at (0,2.297) {$A$};\node[anchor=north east,inner sep=1.5pt] at (0,0) {$F$};\node[anchor=north,inner sep=1.5pt] at (2.04,0) {$B$};\node[rotate=-48.4,anchor=south,inner sep=2pt,font=\footnotesize] at (1.02,1.15) {$384$ m};\node[anchor=north,inner sep=1.5pt,font=\footnotesize] at (1.02,0) {$255$ m};\pfwinkel[0.6][0.62]{(2.04,0)}{131.6}{180}{\beta_1}\end{tikzpicture}}{Eine Seilbahn verbindet die beiden Orte $A$ und $B$. Bestimme die Größe des Winkels $\beta_1$, mit der die Seilbahn von $B$ nach $A$ fährt.}{4}
\par\medskip\noindent\hfill $\beta_1$ $=$\linie{25mm}
\fusshilfe{\gfuss{2}{48,4°}}
\end{pfaufg}
% DRUCK 2022-OS-K5b deckt 2022-OS-K5b
\begin{pfaufg}{nach~P10\\’22}{3.}{}
\gtextkaro{\begin{tikzpicture}[baseline=(current bounding box.north),line width=0.7pt,font=\small]\draw (0,0) -- (7.11,0) -- (4.8,2.96) -- cycle;\draw[dashed] (4.8,2.96) -- (4.8,0);\pfrw[0.25]{(4.8,0)}{90}\node[anchor=north east,inner sep=1.5pt] at (0,0) {$A$};\node[anchor=north west,inner sep=1.5pt] at (7.11,0) {$B$};\node[anchor=south,inner sep=1.5pt] at (4.8,2.96) {$C$};\node[anchor=north,inner sep=1.5pt] at (4.8,0) {$D$};\node[anchor=west,inner sep=1.5pt] at (4.8,1.3) {$h_c$};\node[anchor=south east,inner sep=1.5pt] at (2.4,1.48) {$b = 14{,}1$ m};\pfwinkel[0.9][0.7]{(0,0)}{0}{31.66}{\alpha}\end{tikzpicture}}{Gegeben ist ein Dreieck $ABC$. Das Dreieck ist nicht rechtwinklig. Die Höhe $h_c$ beträgt ca. 7,4~m. Berechne die Größe des Winkels $\alpha$.}{4}
\par\medskip\noindent\hfill $\alpha$ $=$\linie{25mm}
\fusshilfe{\gfuss{3}{31,7°}}
\end{pfaufg}
% DRUCK 2025-OS-K4a deckt 2025-OS-K4a
\begin{pfaufg}{nach~P10\\’25}{4.}{}
\gtextkaro{\begin{tikzpicture}[baseline=(current bounding box.north),line width=0.7pt,font=\small]\fill[black!12] (0,0) -- (6,0) -- (6,0.8) -- cycle;\draw (0,0) -- (6,0) -- (6,0.8) -- cycle;\fill[black!6] (6,0) rectangle (7.4,0.8);\draw (6,0) rectangle (7.4,0.8);\node[anchor=center,inner sep=1.5pt,font=\footnotesize] at (6.7,0.4) {Stufe};\pfrw[0.2]{(6,0)}{90}\node[anchor=east,inner sep=1.5pt,font=\footnotesize] at (5.95,0.45) {$16$ cm};\node[anchor=south,inner sep=1.5pt] at (3,0.4) {$x$};\draw[<->,line width=0.5pt] (0,-0.3) -- (6,-0.3);\node[anchor=north,inner sep=1.5pt,font=\footnotesize] at (3,-0.3) {$170$ cm};\pfwinkel[1.6][0.75]{(0,0)}{0}{7.6}{\alpha}\end{tikzpicture}}{Familie Yücel möchte ihr Haus auch für Personen im Rollstuhl zugänglich machen. An die Stufe der Eingangstür soll eine feste Rampe gebaut werden. Berechne die Länge $x$ der Rampe. Berechne die Größe des Anstiegswinkels $\alpha$.}{6}
\par\medskip\noindent\hfill x $=$\linie{25mm}\quad $\alpha$ $=$\linie{25mm}
\fusshilfe{\gfuss{4}{170,8 cm; 5,4°}}
\end{pfaufg}
\par\medskip\noindent\textcolor{mbgrau}{\rule{\linewidth}{0.3pt}}\par\nobreak\noindent{\footnotesize\textcolor{mbgrau}{2019–2021}}\par\nobreak\smallskip
% DRUCK 2019-OS-K3b deckt 2019-OS-K3b
\begin{pfaufg}{nach~P10\\’19}{5.}{}
\gtextkaro{\begin{tikzpicture}[baseline=(current bounding box.north),line width=0.7pt,font=\small]\draw[dash dot] (0,1.6) -- (0,0) -- (3.22,0) -- cycle;\pfrw[0.25]{(0,0)}{0}\node[anchor=south,inner sep=1.5pt] at (0,1.6) {$A$ (Amir)};\node[anchor=north,inner sep=1.5pt] at (0,0) {$B$ (Ben)};\node[anchor=north,inner sep=1.5pt] at (3.22,0) {$C$ (Chris)};\node[anchor=east,inner sep=1.5pt,font=\footnotesize] at (0,0.8) {$4$ m};\node[anchor=south west,inner sep=1.5pt,font=\footnotesize] at (1.61,0.8) {$9$ m};\pfwinkel[0.55][0.66]{(0,1.6)}{270}{333.6}{\alpha}\end{tikzpicture}}{Ein Fußball-Club hat Training. Drei Kinder spielen sich den Ball zu. Zeige rechnerisch, dass die Größe des Winkels $\alpha$ ungefähr $64^\circ$ beträgt.}{4}
\fusshilfe{\gfuss{5}{63,6°}}
\end{pfaufg}
% DRUCK 2020-OS-K5b deckt 2020-OS-K5b
\begin{pfaufg}{nach~P10\\’20}{6.}{}
\gtextkaro{\begin{tikzpicture}[baseline=(current bounding box.north),line width=0.7pt,font=\small]\draw (0,0) -- (2.8,0) -- (2.8,3) -- (0,1.2) -- cycle;\draw[dashed,line width=0.5pt] (0,1.2) -- (2.8,1.2);\pfrw[0.25]{(0,0)}{0}\pfrw[0.25]{(2.8,0)}{90}\node[anchor=north,inner sep=1.5pt,font=\footnotesize] at (1.4,0) {$14$ cm};\node[anchor=east,inner sep=1.5pt,font=\footnotesize] at (0,0.6) {$6$ cm};\node[anchor=west,inner sep=1.5pt,font=\footnotesize] at (2.8,1.5) {$15$ cm};\pfwinkel[0.55][0.62]{(2.8,3)}{212.7}{270}{\alpha}\end{tikzpicture}}{Eine Schachtel hat die Form eines Prismas. Die Skizze zeigt ihre Seitenfläche. Berechne die Größe des Winkels $\alpha$.}{4}
\par\medskip\noindent\hfill $\alpha$ $=$\linie{25mm}
\fusshilfe{\gfuss{6}{57,3°}}
\end{pfaufg}
% DRUCK 2019-OS-K2d deckt 2019-OS-K2d
\begin{pfaufg}{nach~P10\\’19}{7.}{}
\gtextkaro{\begin{ksys}[xmin=-5,xmax=5,ymin=-5,ymax=5,karo=0.5]\gerade[-1.5]{-2}{2}{f}\gerade[3]{1}{2}{g}\gerade[4]{0}{-2}{h}\punkt{0}{-2}{A}\punkt{2}{-2}{B}\punkt{0}{2}{C}\pfwinkel[1.1][0.62]{(2,-2)}{116.57}{180}{\beta}\end{ksys}}{Im Koordinatensystem bilden die drei Punkte $A$, $B$ und $C$ ein Dreieck. Berechne die Länge der Strecke $\overline{BC}$. Berechne die Größe des Winkels $\beta$.}{6}
\par\medskip\noindent\hfill BC $=$\linie{25mm}\quad $\beta$ $=$\linie{25mm}
\fusshilfe{\gfuss{7}{4,47; 63,4°}}
\end{pfaufg}
\gkopfvor{\gunter{A2}{Seite gesucht}}
% DRUCK 2022-OS-K5d deckt 2022-OS-K5d
\begin{pfaufg}{nach~P10\\’22}{8.}{}
\gtextkaro{\begin{tikzpicture}[baseline=(current bounding box.north),line width=0.7pt,font=\small]\draw (0,0) -- (7.11,0) -- (4.8,2.96) -- cycle;\draw[dashed] (4.8,2.96) -- (4.8,0);\pfrw[0.25]{(4.8,0)}{90}\node[anchor=north east,inner sep=1.5pt] at (0,0) {$A$};\node[anchor=north west,inner sep=1.5pt] at (7.11,0) {$B$};\node[anchor=south,inner sep=1.5pt] at (4.8,2.96) {$C$};\node[anchor=north,inner sep=1.5pt] at (4.8,0) {$D$};\node[anchor=west,inner sep=1.5pt] at (4.8,1.3) {$h_c$};\node[anchor=south west,inner sep=1.5pt] at (5.95,1.48) {$a$};\pfwinkel[0.8][0.62]{(7.11,0)}{128}{180}{52^\circ}\end{tikzpicture}}{Gegeben ist ein Dreieck $ABC$. Das Dreieck ist nicht rechtwinklig. Die Höhe $h_c$ beträgt ca. 7,4~m. Berechne die Länge der Seite $a$.}{4}
\par\medskip\noindent\hfill a $=$\linie{25mm}
\fusshilfe{\gfuss{8}{9,4 m}}
\end{pfaufg}
% DRUCK 2023-OS-K7b deckt 2023-OS-K7b
\begin{pfaufg}{nach~P10\\’23}{9.}{}
\gtextkaro{\begin{tikzpicture}[baseline=(current bounding box.north),line width=0.7pt,font=\small]\fill[black!10] (2.05,0) -- (0,2.05) -- (0,0) -- cycle;\draw (2.05,0) -- (0,2.05) -- (-4.39,0) -- cycle;\draw (0,2.05) -- (0,0);\pfrw[0.25]{(0,0)}{90}\node[anchor=north,inner sep=1.5pt] at (0,0) {$F$};\node[anchor=north west,inner sep=1.5pt] at (2.05,0) {$A$};\node[anchor=south,inner sep=1.5pt] at (0,2.05) {$B$};\node[anchor=north east,inner sep=1.5pt] at (-4.39,0) {$C$};\pfwinkel[0.7][0.66]{(0,2.05)}{205}{270}{65^\circ}\pfwinkel[0.9][0.66]{(0,2.05)}{270}{315}{45^\circ}\pfwinkel[0.6][0.62]{(2.05,0)}{135}{180}{45^\circ}\node[rotate=-45,anchor=south,inner sep=2pt,font=\footnotesize] at (1.025,1.025) {$131{,}5$ cm};\end{tikzpicture}}{Gegeben ist das Dreieck $ABC$. Weise nach, dass die Länge der Strecken $\overline{BF}$ und $\overline{AF}$ jeweils ca. 93,0~cm beträgt. Berechne die Länge der Strecke $\overline{BC}$.}{5}
\par\medskip\noindent\hfill BC $=$\linie{25mm}
\fusshilfe{\gfuss{9}{93,0 cm; 220,0 cm}}
\end{pfaufg}
% DRUCK 2025-OS-K4b deckt 2025-OS-K4b
\begin{pfaufg}{nach~P10\\’25}{10.}{}
\gtextkaro{\begin{tikzpicture}[baseline=(current bounding box.north),line width=0.7pt,font=\small]\fill[black!12] (0,0) -- (6,0) -- (6,0.8) -- cycle;\draw (0,0) -- (6,0) -- (6,0.8) -- cycle;\fill[black!6] (6,0) rectangle (7.4,0.8);\draw (6,0) rectangle (7.4,0.8);\node[anchor=center,inner sep=1.5pt,font=\footnotesize] at (6.7,0.4) {Stufe};\pfrw[0.2]{(6,0)}{90}\node[anchor=east,inner sep=1.5pt,font=\footnotesize] at (5.95,0.45) {$16$ cm};\node[anchor=south,inner sep=1.5pt] at (3,0.4) {$x$};\draw[<->,line width=0.5pt] (0,-0.3) -- (6,-0.3);\node[anchor=north,inner sep=1.5pt,font=\footnotesize] at (3,-0.3) {$170$ cm};\pfwinkel[1.6][0.75]{(0,0)}{0}{7.6}{\alpha}\end{tikzpicture}}{Familie Yücel möchte ihr Haus auch für Personen im Rollstuhl zugänglich machen. An die Stufe der Eingangstür soll eine feste Rampe gebaut werden. Damit eine Person im Rollstuhl eine Rampe ohne Hilfe hochfahren kann, darf die Steigung einer Rampe maximal 6~\% betragen. Vervollständige den folgenden Satz: Eine Steigung von 6~\% bedeutet, dass auf einer waagerechten Strecke von \linie{15mm} Metern ein Höhenunterschied von \linie{15mm} Metern erreicht wird. Frau Yücel behauptet: „Unsere Rampe wird zu steil, da die Steigung mehr als 6~\% beträgt.“ Hat Frau Yücel recht? Überprüfe die Aussage.}{4}
\fusshilfe{\gfuss{10}{100 und 6; ja}}
\end{pfaufg}
% DRUCK 2026-FOR-K4c deckt 2026-FOR-K4c
\begin{pfaufg}{nach~P10\\’26}{11.\textsuperscript{*}}{}
\gtextkaro{\begin{tikzpicture}[baseline=(current bounding box.north),line width=0.7pt,font=\small]\draw (0,0) -- (4.6,0) -- (5.9,2.92) -- (1.3,2.92) -- cycle;\draw[densely dotted] (4.6,0) -- (5.9,0) -- (5.9,2.92);\pfrw[0.25]{(5.9,0)}{90}\node[anchor=north east,inner sep=1.5pt] at (0,0) {$A$};\node[anchor=north east,inner sep=1.5pt] at (4.6,0) {$B$};\node[anchor=north west,inner sep=1.5pt] at (5.9,0) {$F$};\node[anchor=south west,inner sep=1.5pt] at (5.9,2.92) {$C$};\node[anchor=south east,inner sep=1.5pt] at (1.3,2.92) {$D$};\node[anchor=north,inner sep=1.5pt,font=\footnotesize] at (5.25,0) {$13$ cm};\node[anchor=west,inner sep=1.5pt] at (5.9,1.46) {$h$};\node[rotate=66,anchor=south,inner sep=2pt,font=\footnotesize] at (5.25,1.46) {$32$ cm};\draw (0,0) -- (5.9,2.92);\pfwinkel[1.6][0.75]{(0,0)}{0}{26.3}{22^\circ}\end{tikzpicture}}{Bei dem Viereck $ABCD$ sind die gegenüberliegenden Seiten jeweils parallel zueinander. Bestimme die Länge der Diagonale $\overline{AC}$.}{5}
\par\medskip\noindent\hfill AC $=$\linie{25mm}
\fusshilfe{\gfuss{11}{78,1 cm}}
\end{pfaufg}
\par\medskip\noindent\textcolor{mbgrau}{\rule{\linewidth}{0.3pt}}\par\nobreak\noindent{\footnotesize\textcolor{mbgrau}{2016–2021}}\par\nobreak\smallskip
% DRUCK 2021-OS-K3a deckt 2021-OS-K3a
% DRUCK 2021-OS-K3b deckt 2021-OS-K3b
\begin{pfaufg}{nach~P10\\’21}{12.}{}
\gtextkaro{\begin{tikzpicture}[baseline=(current bounding box.north),line width=0.7pt,font=\small]\draw (0,0) -- (0,1.966) -- (-2.732,2.906) -- (-1.377,0) -- cycle;\draw (-1.377,0) -- (0,1.966);\pfrw[0.25]{(0,0)}{90}\node[anchor=north west,inner sep=1.5pt] at (0,0) {$A$};\node[anchor=south west,inner sep=1.5pt] at (0,1.966) {$B$};\node[anchor=south east,inner sep=1.5pt] at (-2.732,2.906) {$C$};\node[anchor=north east,inner sep=1.5pt] at (-1.377,0) {$D$};\node[rotate=55,anchor=north,inner sep=2pt,font=\footnotesize] at (-0.69,0.98) {$1500$ m};\pfwinkel[0.55][0.62]{(-1.377,0)}{0}{55}{55^\circ}\end{tikzpicture}}{Gegeben ist das Viereck $ABCD$.\begin{gbteil}{a)}{}{}Zeige mit einer Rechnung, dass die Länge der Strecke $\overline{AB}$ ca. 1229~m beträgt.\end{gbteil}}{4}
\begin{gbteil}{b)}{}{}
Berechne die Länge der Strecke $\overline{AD}$.
\karo{4}
\par\medskip\noindent\hfill AD $=$\linie{25mm}
\end{gbteil}
\fusshilfe{\gfuss{12}{a) 1228,7 m; b) 860,4 m}}
\end{pfaufg}
% DRUCK 2017-OS-K4b deckt 2017-OS-K4b
\begin{pfaufg}{nach~P10\\’17}{13.}{}
\gtextkaro{\begin{tikzpicture}[baseline=(current bounding box.north),line width=0.7pt,font=\small]\draw (0,0) -- (4.68,0) -- (1.235,2.401) -- cycle;\node[anchor=north east,inner sep=1.5pt] at (0,0) {$A$};\node[anchor=north west,inner sep=1.5pt] at (4.68,0) {$B$};\node[anchor=south,inner sep=1.5pt] at (1.235,2.401) {$C$};\pfwinkel[0.8][0.68]{(4.68,0)}{145.1}{180}{34{,}9^\circ}\draw[line width=1.6pt] (3.648,0) -- (3.648,0.72);\node[anchor=north,inner sep=1.5pt] at (3.648,0) {$P$};\node[anchor=east,inner sep=1.5pt,font=\footnotesize] at (3.6,0.36) {$1{,}20$ m};\end{tikzpicture}}{Der Dachboden eines Hauses soll ausgebaut werden. Die Skizze zeigt den Querschnitt des Dachbodens. Im Dachboden soll eine senkrechte Wand gebaut werden. Sie reicht bis zur Dachschräge $\overline{BC}$ und ist 1,20~m hoch. Berechne, wie viel Meter der Fußpunkt $P$ der Wand vom Punkt $B$ entfernt sein muss.}{4}
\par\medskip\noindent\hfill PB $=$\linie{25mm}
\fusshilfe{\gfuss{13}{1,72 m}}
\end{pfaufg}
% DRUCK 2018-OS-K4d deckt 2018-OS-K4d
\begin{pfaufg}{nach~P10\\’18}{14.}{}
\gtextkaro{\begin{tikzpicture}[baseline=(current bounding box.north),line width=0.7pt,font=\small]\draw[dashed,line width=0.5pt] (-0.3,0) -- (5.2,0);\draw[line width=1.2pt] (0,0) -- (0,0.5);\draw (0,0.5) -- (4.6,0.5);\draw[dashed] (4.6,0.5) -- (4.6,3.156);\draw[line width=1.6pt] (4.6,3.156) -- (4.6,3.721);\draw (0,0.5) -- (4.6,3.156);\draw (0,0.5) -- (4.6,3.721);\pfrw[0.25]{(4.6,0.5)}{90}\node[anchor=south east,inner sep=1.5pt] at (0,0.5) {$A$};\node[anchor=north west,inner sep=1.5pt] at (4.6,0.5) {$B$};\node[anchor=west,inner sep=1.5pt] at (4.6,3.156) {$C$};\node[anchor=south,inner sep=1.5pt] at (4.6,3.721) {$D$};\node[anchor=west,inner sep=1.5pt,font=\footnotesize] at (4.68,3.44) {$12{,}0$ m};\node[anchor=east,inner sep=1.5pt,font=\footnotesize] at (-0.05,0.25) {$1{,}5$ m};\pfwinkel[1.3][0.75]{(0,0.5)}{0}{30}{30^\circ}\node[rotate=30,anchor=north,inner sep=2pt,font=\footnotesize] at (2.9,2.17) {$112{,}8$ m};\end{tikzpicture}}{Tim möchte mit einem Messgerät die Höhe des Hasenberges bestimmen. Auf der Bergspitze $C$ befindet sich eine 12,0~m hohe Aussichtsplattform $D$. Vom Punkt $A$ misst er die zwei Winkelgrößen $30^\circ$ und $35^\circ$. Der Punkt $A$ liegt 1,5~m über dem Boden. Tim hat für $\overline{AC}$ eine Länge von 112,8~m bestimmt. Ermittle die Höhe des Hasenberges. Gib an, auf welcher Höhe über dem Boden sich die Aussichtsplattform $D$ befindet.}{5}
\par\medskip\noindent\hfill Berg $=$\linie{25mm}\quad D $=$\linie{25mm}
\fusshilfe{\gfuss{14}{57,9 m; 69,9 m}}
\end{pfaufg}
% DRUCK 2016-OS-K7c deckt 2016-OS-K7c
\begin{pfaufg}{nach~P10\\’16}{15.\textsuperscript{*}}{}
\gtextkaro{\begin{tikzpicture}[baseline=(current bounding box.north),line width=0.7pt,font=\small]\draw[mbgrau,line width=0.5pt] (-0.4,0.18) -- (4.8,0.18);\node[anchor=south,inner sep=1.5pt,font=\footnotesize,text=mbgrau] at (1.1,0.18) {Spree};\draw[line width=1.1pt] (2.164,-2.977) -- (0,0) -- (2.164,0);\draw[dashed,line width=1.1pt] (2.164,0) -- (4.172,0);\draw (2.164,-2.977) -- (4.172,0);\draw[line width=0.5pt] (2.164,-2.977) -- (2.164,0);\pfrw[0.25]{(2.164,0)}{270}\node[anchor=south east,inner sep=1.5pt] at (0,0) {$C$};\node[anchor=south,inner sep=1.5pt] at (2.164,0) {$D$};\node[anchor=south west,inner sep=1.5pt] at (4.172,0) {$B$};\node[anchor=north,inner sep=1.5pt] at (2.164,-2.977) {$A$};\node[anchor=north east,inner sep=1.5pt,font=\footnotesize] at (-0.05,-0.05) {Bioladen};\node[anchor=south,inner sep=1.5pt,font=\footnotesize] at (2.164,0.45) {Brücke};\node[anchor=north west,inner sep=1.5pt,font=\footnotesize] at (4.2,-0.1) {Anlegestelle};\node[anchor=north,inner sep=1.5pt,font=\footnotesize] at (2.164,-3.25) {Ausflugslokal};\node[anchor=south,inner sep=1.5pt,font=\footnotesize] at (1.08,0) {$541$ m};\node[anchor=north east,inner sep=1.5pt,font=\footnotesize] at (0.9,-1.6) {$920$ m};\pfwinkel[0.6][0.62]{(0,0)}{306}{360}{54^\circ}\pfwinkel[0.7][0.62]{(2.164,-2.977)}{56}{126}{70^\circ}\end{tikzpicture}}{Im Spreewald gibt es viele sumpfige Wiesen. Um vom Ausflugslokal $A$ zur Brücke bei $D$ zu gelangen, musste man bisher am Bioladen $C$ vorbei. Um von der Anlegestelle $B$ direkt zur Brücke bei $D$ gehen zu können, wird ein neuer Weg $\overline{BD}$ gebaut. Ermittle die Länge des Weges $\overline{BD}$.}{5}
\par\medskip\noindent\hfill BD $=$\linie{25mm}
\fusshilfe{\gfuss{15}{502 m}}
\end{pfaufg}

% ===================== B Berechnungen in beliebigen Dreiecken =====================
\gkopfvor{\gbstvor{B}\pfstufekopf{Berechnungen in beliebigen Dreiecken}{}}
% DRUCK 2025-OS-K4c deckt 2025-OS-K4c
\begin{pfaufg}{nach~P10\\’25}{16.\textsuperscript{*}}{}
\gtextkaro{\begin{tikzpicture}[baseline=(current bounding box.north),line width=0.7pt,font=\small]\fill[black!15] (3.2,0) -- (3.2,0.87) -- (5.344,0.87) -- (5.344,1.7365) -- (6.6,1.7365) -- (6.6,0) -- cycle;\draw (0,0) -- (6.6,0);\draw (3.2,0) -- (3.2,0.87) -- (5.344,0.87) -- (5.344,1.7365) -- (6.6,1.7365);\draw[line width=1.3pt] (0,0) -- (5.344,1.7365);\draw[dashed,line width=0.5pt] (3.2,0) -- (5.344,1.7365);\node[anchor=south,inner sep=1.5pt] at (2.4,0.85) {$y$};\pfwinkel[1.7][0.8]{(0,0)}{0}{18}{5^\circ}\pfwinkel[0.35][0.75]{(3.2,0)}{39}{180}{}\node[anchor=center,inner sep=1.5pt,font=\footnotesize] at (2.8,0.3) {$141^\circ$};\draw[<->,line width=0.5pt] (0,-0.3) -- (3.2,-0.3);\node[anchor=north,inner sep=1.5pt,font=\footnotesize] at (1.6,-0.3) {$160$ cm};\end{tikzpicture}}{Für die Treppe am Hinterausgang müssen zwei Stufen überwunden werden. Dafür wird eine mobile Rampe an die Treppe angelegt. Berechne die Länge $y$ der Rampe.}{4}
\par\medskip\noindent\hfill y $=$\linie{25mm}
\fusshilfe{\gfuss{16}{180,1 cm}}
\end{pfaufg}\vspace{-6pt}
% DRUCK 2024-OS-K6d deckt 2024-OS-K6d
\begin{pfaufg}{nach~P10\\’24}{17.\textsuperscript{*}}{}
\gtextkaro{\begin{tikzpicture}[baseline=(current bounding box.north),line width=0.7pt,font=\small]\draw (0,2.297) -- (0,0) -- (2.04,0) -- cycle;\pfrw[0.25]{(0,0)}{0}\node[anchor=south,inner sep=1.5pt] at (0,2.297) {$A$};\node[anchor=north east,inner sep=1.5pt] at (0,0) {$F$};\node[anchor=north,inner sep=1.5pt] at (2.04,0) {$B$};\node[rotate=-48.4,anchor=south,inner sep=2pt,font=\footnotesize] at (1.02,1.15) {$384$ m};\draw (2.04,0) -- (5.139,1.354) -- (0,2.297);\node[anchor=south west,inner sep=1.5pt] at (5.139,1.354) {$C$};\pfwinkel[0.9][0.72]{(0,2.297)}{311.6}{349.6}{\alpha}\pfwinkel[0.45][0.7]{(2.04,0)}{23.6}{131.6}{\beta_2}\end{tikzpicture}}{Eine Seilbahn verbindet die beiden Orte $A$ und $B$. Der Ort $B$ und der Ort $C$ sind mit einer weiteren Seilbahn verbunden. Folgende Angaben sind bekannt: $\alpha = 38^\circ$ und $\beta_2 = 108^\circ$. Berechne die Länge der Seilbahnstrecke von $B$ nach $C$.}{4}
\par\medskip\noindent\hfill BC $=$\linie{25mm}
\fusshilfe{\gfuss{17}{422,8 m}}
\end{pfaufg}\vspace{-6pt}
\par\medskip\noindent\textcolor{mbgrau}{\rule{\linewidth}{0.3pt}}\par\nobreak\noindent{\footnotesize\textcolor{mbgrau}{2014–2021}}\par\nobreak\smallskip
% DRUCK 2015-OS-K5d deckt 2015-OS-K5d
\begin{pfaufg}{nach~P10\\’15}{18.}{}
\gtextkaro{\begin{tikzpicture}[baseline=(current bounding box.north),line width=0.7pt,font=\small]\draw (0,0) -- (3.445,-1.322) -- (5.265,0) -- (3.445,1.322) -- cycle;\draw (0,0) -- (5.265,0);\draw (3.445,1.322) -- (3.445,-1.322);\pfrw[0.2]{(3.445,0)}{0}\node[anchor=east,inner sep=1.5pt] at (0,0) {$D$};\node[anchor=west,inner sep=1.5pt] at (5.265,0) {$B$};\node[anchor=north,inner sep=1.5pt] at (3.445,-1.322) {$A$};\node[anchor=south,inner sep=1.5pt] at (3.445,1.322) {$C$};\node[anchor=south east,inner sep=1.5pt,font=\footnotesize] at (3.445,0) {$E$};\node[anchor=north east,inner sep=1.5pt,font=\footnotesize] at (1.72,-0.66) {$4{,}1$ cm};\pfwinkel[0.45][0.0]{(3.445,-1.322)}{36}{159}{}\node[anchor=south,inner sep=1.5pt,font=\scriptsize] at (3.445,-0.85) {$123^\circ$};\pfwinkel[0.8][0.72]{(5.265,0)}{180}{216}{36^\circ}\pfwinkel[1.3][0.62]{(0,0)}{-21}{0}{\delta}\end{tikzpicture}}{Gegeben ist das Drachenviereck $ABCD$; die Diagonalen schneiden sich in $E$. Der Flächeninhalt des Dreiecks $AED$ soll ermittelt werden. Gib die benötigten Seiten und Winkel an. Notiere eine Reihenfolge der Schritte, die zur Berechnung des gesuchten Flächeninhaltes notwendig sind.}{6}
\fusshilfe{\gfuss{18}{21° und rechter Winkel bei E; 2,81 cm²}}
\end{pfaufg}\vspace{-6pt}
% DRUCK 2018-OS-K4c deckt 2018-OS-K4c
\begin{pfaufg}{nach~P10\\’18}{19.\textsuperscript{*}}{}
\gtextkaro{\begin{tikzpicture}[baseline=(current bounding box.north),line width=0.7pt,font=\small]\draw[dashed,line width=0.5pt] (-0.3,0) -- (5.2,0);\draw[line width=1.2pt] (0,0) -- (0,0.5);\draw (0,0.5) -- (4.6,0.5);\draw[dashed] (4.6,0.5) -- (4.6,3.156);\draw[line width=1.6pt] (4.6,3.156) -- (4.6,3.721);\draw (0,0.5) -- (4.6,3.156);\draw (0,0.5) -- (4.6,3.721);\pfrw[0.25]{(4.6,0.5)}{90}\node[anchor=south east,inner sep=1.5pt] at (0,0.5) {$A$};\node[anchor=north west,inner sep=1.5pt] at (4.6,0.5) {$B$};\node[anchor=west,inner sep=1.5pt] at (4.6,3.156) {$C$};\node[anchor=south,inner sep=1.5pt] at (4.6,3.721) {$D$};\node[anchor=west,inner sep=1.5pt,font=\footnotesize] at (4.68,3.44) {$12{,}0$ m};\node[anchor=east,inner sep=1.5pt,font=\footnotesize] at (-0.05,0.25) {$1{,}5$ m};\pfwinkel[1.3][0.75]{(0,0.5)}{0}{30}{30^\circ}\pfwinkel[2.1][0.88]{(0,0.5)}{30}{35}{\gamma}\pfwinkel[0.5][0.62]{(4.6,3.721)}{215}{270}{\delta}\pfwinkel*[2.6]{(0,0.5)}{0}{35}{}\node[anchor=south,inner sep=1.5pt,font=\footnotesize] at (2.75,0.55) {$35^\circ$};\end{tikzpicture}}{Tim möchte mit einem Messgerät die Höhe des Hasenberges bestimmen. Auf der Bergspitze $C$ befindet sich eine 12,0~m hohe Aussichtsplattform $D$. Vom Punkt $A$ misst er die zwei Winkelgrößen $30^\circ$ und $35^\circ$. Der Punkt $A$ liegt 1,5~m über dem Boden. Tim gibt die Größe des Winkels $\delta$ mit $55^\circ$ an. Er hat für $\overline{AC}$ eine Länge von 112,8~m bestimmt. Bestätige diese Angabe durch eine geeignete Rechnung.}{4}
\fusshilfe{\gfuss{19}{112,8 m}}
\end{pfaufg}\vspace{-6pt}
% DRUCK 2015-OS-K5c deckt 2015-OS-K5c
\begin{pfaufg}{nach~P10\\’15}{20.\textsuperscript{*}}{}
\gtextkaro{\begin{tikzpicture}[baseline=(current bounding box.north),line width=0.7pt,font=\small]\draw (0,0) -- (3.445,-1.322) -- (5.265,0) -- (3.445,1.322) -- cycle;\draw (0,0) -- (5.265,0);\draw (3.445,1.322) -- (3.445,-1.322);\pfrw[0.2]{(3.445,0)}{0}\node[anchor=east,inner sep=1.5pt] at (0,0) {$D$};\node[anchor=west,inner sep=1.5pt] at (5.265,0) {$B$};\node[anchor=north,inner sep=1.5pt] at (3.445,-1.322) {$A$};\node[anchor=south,inner sep=1.5pt] at (3.445,1.322) {$C$};\node[anchor=south east,inner sep=1.5pt,font=\footnotesize] at (3.445,0) {$E$};\node[anchor=north east,inner sep=1.5pt,font=\footnotesize] at (1.72,-0.66) {$4{,}1$ cm};\pfwinkel[0.45][0.0]{(3.445,-1.322)}{36}{159}{}\node[anchor=south,inner sep=1.5pt,font=\scriptsize] at (3.445,-0.85) {$123^\circ$};\pfwinkel[0.8][0.72]{(5.265,0)}{180}{216}{36^\circ}\end{tikzpicture}}{Gegeben ist das Drachenviereck $ABCD$. Berechne die Länge der Strecke $\overline{BD}$.}{4}
\par\medskip\noindent\hfill BD $=$\linie{25mm}
\fusshilfe{\gfuss{20}{5,85 cm}}
\end{pfaufg}\vspace{-6pt}
% DRUCK 2017-OS-K4c deckt 2017-OS-K4c
\begin{pfaufg}{nach~P10\\’17}{21.\textsuperscript{*}}{}
\gtextkaro{\begin{tikzpicture}[baseline=(current bounding box.north),line width=0.7pt,font=\small]\draw (0,0) -- (4.68,0) -- (1.235,2.401) -- cycle;\node[anchor=north east,inner sep=1.5pt] at (0,0) {$A$};\node[anchor=north west,inner sep=1.5pt] at (4.68,0) {$B$};\node[anchor=south,inner sep=1.5pt] at (1.235,2.401) {$C$};\pfwinkel[0.8][0.68]{(4.68,0)}{145.1}{180}{34{,}9^\circ}\pfwinkel[0.6][0.62]{(0,0)}{0}{62.8}{62{,}8^\circ}\node[rotate=62.8,anchor=south,inner sep=2pt,font=\footnotesize] at (0.62,1.2) {$4{,}50$ m};\fill[black!18] (5.40,0.00) -- (5.71,0.60) -- (6.91,1.30) -- (6.60,0.70) -- cycle;\fill[black!28] (6.57,0.00) -- (5.71,0.60) -- (6.91,1.30) -- (7.77,0.70) -- cycle;\draw[line width=0.5pt] (5.40,0.00) -- (6.57,0.00) -- (5.71,0.60) -- cycle;\draw[line width=0.5pt] (7.77,0.70) -- (6.91,1.30) -- (6.60,0.70);\draw[line width=0.5pt] (5.40,0.00) -- (6.60,0.70) (6.57,0.00) -- (7.77,0.70) (5.71,0.60) -- (6.91,1.30);\draw[line width=0.5pt,densely dashed] (6.60,0.70) -- (7.77,0.70);\node[anchor=north,inner sep=1.5pt,font=\footnotesize] at (5.40,0.00) {$A$};\node[anchor=north,inner sep=1.5pt,font=\footnotesize] at (6.57,0.00) {$B$};\node[anchor=south east,inner sep=1.5pt,font=\footnotesize] at (5.71,0.60) {$C$};\node[anchor=west,inner sep=1.5pt,font=\footnotesize] at (7.2,0.3) {$8{,}00$ m};\end{tikzpicture}}{Der Dachboden eines Hauses soll ausgebaut werden. Die Skizze zeigt den Querschnitt des Dachbodens. Zur Wärmedämmung sollen die beiden gefärbten Dachflächen isoliert werden. Berechne die Länge der Dachschräge $\overline{BC}$. Der Dachboden hat eine Länge von 8,00~m. Berechne, wie viel $\mathrm{m^2}$ Dachfläche isoliert werden müssen.}{6}
\par\medskip\noindent\hfill BC $=$\linie{25mm}\quad A $=$\linie{25mm}
\fusshilfe{\gfuss{21}{7,00 m; 92,0 m²}}
\end{pfaufg}\vspace{-6pt}
% DRUCK 2014-OS-K2b deckt 2014-OS-K2b
\begin{pfaufg}{nach~P10\\’14}{22.\textsuperscript{*}}{}
\gtextkaro{\begin{tikzpicture}[baseline=(current bounding box.north),line width=0.7pt,font=\small]\draw (0,0) -- (2.68,2.25) -- (4.035,-1.469) -- (1.51,-0.55) -- cycle;\fill (0,0) circle (1.3pt);\fill (2.68,2.25) circle (1.3pt);\fill (4.035,-1.469) circle (1.3pt);\fill (1.51,-0.55) circle (1.3pt);\node[anchor=east,inner sep=1.5pt,font=\footnotesize] at (0,0) {Forsthaus};\node[anchor=south,inner sep=1.5pt,font=\footnotesize] at (2.68,2.25) {Jugendherberge};\node[anchor=north,inner sep=1.5pt,font=\footnotesize] at (4.035,-1.469) {Restaurant};\node[anchor=north,inner sep=1.5pt,font=\footnotesize] at (1.51,-0.55) {Badestelle};\node[rotate=40,anchor=south,inner sep=2pt,font=\footnotesize] at (1.34,1.12) {$2800$ m};\node[anchor=north east,inner sep=1.5pt,font=\footnotesize] at (0.75,-0.27) {$1{,}5$ km};\node[anchor=north east,inner sep=1.5pt,font=\footnotesize] at (2.77,-1.01) {$2500$ m};\pfwinkel[0.7][0.68]{(0,0)}{-20}{40}{60^\circ}\pfwinkel[0.7][0.68]{(4.035,-1.469)}{110}{160}{50^\circ}\end{tikzpicture}}{Zwei Wandergruppen aus Berlin und Potsdam sind in einer Jugendherberge untergebracht. Beide Gruppen wollen zur Badestelle laufen. Die Berliner Gruppe wandert zuerst zum Restaurant und danach zur Badestelle. Berechne die Gesamtlänge dieses Weges.}{5}
\fusshilfe{\gfuss{22}{5665 m}}
\end{pfaufg}\vspace{-6pt}
% DRUCK 2021-OS-K3c deckt 2021-OS-K3c
\begin{pfaufg}{nach~P10\\’21}{23.\textsuperscript{*}}{}
\gtextkaro{\begin{tikzpicture}[baseline=(current bounding box.north),line width=0.7pt,font=\small]\draw (0,0) -- (0,1.966) -- (-2.732,2.906) -- (-1.377,0) -- cycle;\draw (-1.377,0) -- (0,1.966);\pfrw[0.25]{(0,0)}{90}\node[anchor=north west,inner sep=1.5pt] at (0,0) {$A$};\node[anchor=south west,inner sep=1.5pt] at (0,1.966) {$B$};\node[anchor=south east,inner sep=1.5pt] at (-2.732,2.906) {$C$};\node[anchor=north east,inner sep=1.5pt] at (-1.377,0) {$D$};\node[rotate=55,anchor=north,inner sep=2pt,font=\footnotesize] at (-0.69,0.98) {$1500$ m};\pfwinkel[0.55][0.62]{(-1.377,0)}{0}{55}{55^\circ}\node[rotate=-65,anchor=south,inner sep=2pt,font=\footnotesize] at (-2.05,1.45) {$2004$ m};\pfwinkel[0.85][0.72]{(-1.377,0)}{55}{115}{60^\circ}\pfwinkel[0.45][0.72]{(0,1.966)}{161}{270}{109^\circ}\end{tikzpicture}}{Gegeben ist das Viereck $ABCD$. Prüfe, ob folgende Aussage stimmt: „Die Strecke $\overline{BC}$ ist genauso lang wie die Strecke $\overline{DC}$.“ Entscheide mit Hilfe einer Rechnung.}{6}
\fusshilfe{\gfuss{23}{nein}}
\end{pfaufg}\vspace{-6pt}
% DRUCK 2020-OS-K7c deckt 2020-OS-K7c
\begin{pfaufg}{nach~P10\\’20}{24.\textsuperscript{*}}{}
\gtextkaro{\begin{tikzpicture}[baseline=(current bounding box.north),line width=0.7pt,font=\small]\draw (0,0) -- (5.359,1.436) -- (0.736,1.51) -- cycle;\fill (1.302,0.349) circle (1.3pt);\node[anchor=north east,inner sep=1.5pt] at (0,0) {$D$};\node[anchor=west,inner sep=1.5pt] at (5.359,1.436) {$E$};\node[anchor=south,inner sep=1.5pt] at (0.736,1.51) {$F$};\node[anchor=north,inner sep=1.5pt] at (1.302,0.349) {$P$};\node[anchor=east,inner sep=1.5pt,font=\footnotesize] at (0.368,0.755) {$4$ m};\node[anchor=south,inner sep=1.5pt,font=\footnotesize] at (3.05,1.47) {$11$ m};\pfwinkel[0.4][0.0]{(0.736,1.51)}{244}{359.1}{}\node[anchor=center,inner sep=1.5pt,font=\scriptsize] at (1.05,1.1) {$115^\circ$};\pfwinkel[1.0][0.72]{(5.359,1.436)}{179.1}{195}{16^\circ}\end{tikzpicture}}{Betrachtet wird das Dreieck $DEF$. Der Punkt $P$ liegt auf der Strecke $\overline{DE}$. Die Länge der Strecke $\overline{PE}$ beträgt 10~m. Berechne die Länge der Strecke $\overline{DP}$.}{5}
\par\medskip\noindent\hfill DP $=$\linie{25mm}
\fusshilfe{\gfuss{24}{3,2 m}}
\end{pfaufg}\vspace{-6pt}
% DRUCK 2019-OS-K3c deckt 2019-OS-K3c
\begin{pfaufg}{nach~P10\\’19}{25.\textsuperscript{*}}{}
\gtextkaro{\begin{tikzpicture}[baseline=(current bounding box.north),line width=0.7pt,font=\small]\draw[dash dot] (0,1.775) -- (0,0) -- (3.64,0) -- (3.456,2.637) -- cycle;\draw[dashed] (0,1.775) -- (3.64,0);\pfrw[0.25]{(0,0)}{0}\node[anchor=east,inner sep=1.5pt] at (0,1.775) {Amir $A$};\node[anchor=north,inner sep=1.5pt] at (0,0) {Ben $B$};\node[anchor=north,inner sep=1.5pt] at (3.64,0) {$C$ Chris};\node[anchor=south,inner sep=1.5pt] at (3.456,2.637) {$D$ Deniz};\node[anchor=south west,inner sep=1.5pt,font=\footnotesize] at (1.82,0.89) {$9$ m};\pfwinkel[0.55][0.66]{(0,1.775)}{270}{334}{64^\circ}\pfwinkel[0.5][0.66]{(3.64,0)}{94}{180}{86^\circ}\pfwinkel[0.5][0.66]{(3.456,2.637)}{194}{274}{80^\circ}\end{tikzpicture}}{Beim Fußballtraining trainieren vier Kinder. Die Bälle werden in einem bestimmten Winkel angenommen und weitergespielt. Berechne den Abstand zwischen Amir und Deniz.}{5}
\par\medskip\noindent\hfill AD $=$\linie{25mm}
\fusshilfe{\gfuss{25}{7,9 m}}
\end{pfaufg}\vspace{-6pt}
\end{document}